\documentclass{article}
\usepackage{listings}
\usepackage{datetime}
\usepackage[margin=1in]{geometry}

\lstset{
    basicstyle=\small\ttfamily,
    breaklines=true,
    frame=single,
    tabsize=4,
    showstringspaces=false
}

\title{HW5: SQL Part 1 - Queries}
\author{Abdullah Junejo}
\date{\today}

\begin{document}

\maketitle

\section*{Creating Table for EmployeeAttrition1.csv}

\begin{lstlisting}[language=SQL]
CREATE TABLE EmployeeAttrition1 (
    Employee_Number INT, Age INT, Business_Travel varchar(255), Daily_Rate INT,
    Department varchar(255), Distance_from_Home INT, Education INT,
    Education_Field varchar(255), Environment_Satisfaction INT, 
    Gender varchar(255), Hourly_Rate INT, Job_Involvement INT, 
    Job_Level INT, Job_Role varchar(255), Job_Satisfaction INT,
    Marital_Status varchar(255), Monthly_Income INT, Monthly_Rate INT,
    Num_Companies_Worked INT, Percent_Salary_Hike INT, 
    Performance_Rating INT, Relationship_Satisfaction INT,
    Standard_Hours INT, Stock_Option_Level INT, Total_Working_Years INT, 
    Training_Times_Last_Year INT, Work_Life_Balance INT,
    Years_At_Company INT, Years_In_CurrentRole INT, 
    Years_Since_Last_Promotion INT, Years_With_Curr_Manager INT
);
\end{lstlisting}

\section*{Creating Table for EmployeeAttrition2.csv}

\begin{lstlisting}[language=SQL]
CREATE TABLE EmployeeAttritionTable2 (
    Employee_Number INT, Over_18 varchar(64), Over_Time varchar(64),
    Attrition varchar(64)
);
\end{lstlisting}

\section*{Queries}

\subsection*{Q1}
\begin{lstlisting}[language=SQL]
SELECT count(Employee_Number) as Total_Records 
FROM EmployeeAttrition1;
\end{lstlisting}

\subsection*{Q2}
\begin{lstlisting}[language=SQL]
SELECT Job_Role as JobRole, count(Job_Role) as Count_Job_Role 
FROM EmployeeAttrition1 
GROUP BY Job_Role 
ORDER BY Count_Job_Role DESC;
\end{lstlisting}

\subsection*{Q3}
\begin{lstlisting}[language=SQL]
SELECT Job_Role as JobRolee, 
       avg(Monthly_Income) as Average_Monthly_Income, 
       avg(Percent_Salary_Hike) as Salary_Hike_Increase 
FROM EmployeeAttrition1 
GROUP BY Job_Role 
ORDER BY JobRolee;
\end{lstlisting}

\subsection*{Q4}
\begin{lstlisting}[language=SQL]
SELECT avg(Job_Satisfaction) as Avg_Job_Satisfaction, 
       Gender, Marital_Status as Current_Gender_Marital_Status 
FROM EmployeeAttrition1 
GROUP BY Gender, Marital_Status;
\end{lstlisting}

\subsection*{Q5}
\begin{lstlisting}[language=SQL]
SELECT Job_Role as Role, min(Age) as Min_Age, 
       max(Age) as Max_Age, min(Hourly_Rate) as Min_Hourly_Rate, 
       max(Hourly_Rate) as Max_Hourly_Rate 
FROM EmployeeAttrition1 
GROUP BY Job_Role;
\end{lstlisting}

\subsection*{Q6}
\begin{lstlisting}[language=SQL]
SELECT Emp1.Employee_Number as Emp_Number, Emp1.Age, 
       Emp1.Gender, Emp1.Job_Role as Role, 
       Emp2.Over_Time as OverTimePeriod, Emp2.Attrition 
FROM EmployeeAttrition1 Emp1 
JOIN EmployeeAttritionTable2 Emp2 
ON Emp1.Employee_Number = Emp2.Employee_Number 
LIMIT 20;
\end{lstlisting}

\end{document}
